%!TEX root = ./poly.tex

\section{Norms}
\label{sec:norms}

In this section we study the Norm map in $\RO(G)$-graded degrees. We briefly study $p$-primary Adams operations in \sec\ref{sub:adams}. We determine Norms for the de Rham-Witt complex in \sec\ref{sub:dRW-norm}.

\subsection{\tops{$p$}{p}-primary Adams operations}
\label{sub:adams}

The identity $FNf = f^p$ nearly determines the norm in ghost coordinates; to fully determine it, we need to understand the composite $R^n N$. This will turn out to be the $p^\th$ Adams operation $\psi^p$, so in this subsection we recall this operation and prove some facts about it.

\begin{definition}
  For an integer $\ell$, the natural map
  \[ \psi^\ell\colon \THH(R) = R^{\tnsr\T} \to R^{\tnsr\T/C_\ell} = \THH(R)_{\tnsr C_\ell} \]
  is called the \emph{$\ell^\th$ Adams operation}.
\end{definition}

Beware that $\THH(R)$ is equivalent to $\THH(R)_{\tnsr C_\ell}$ in $\Sp$, but not in $\Sp^{B\T}$. Therefore, $\psi^\ell$ may be regarded as an endomorphism of $\THH$, but it does not necessarily pass to $\TC^-$ and $\TP$. However, $\THH(R)$ and $\THH(R)_{\tnsr C_\ell}$ do become equivalent in $\Sp[1/\ell]^{B\T}$, so that in this case $\psi^\ell$ may be regarded as an automorphism of $\TC^-$ and $\TP$. This holds in particular if we $p$-complete at a prime $p$ not dividing $\ell$ (in which case $\ell$ may be generalized to any element of $\Z_p^\times$).

  For this reason, the analysis in the literature has mostly focused on the case $p\nmid\ell$. For example, in this case $\psi^\ell$ acts as multiplication by $\ell^i$ on $\gr^i$ of $\THH(-;\Z_p)$, $\TCmin(-;\Z_p)$, and $\TP(-;\Z_p)$ by \cite[Proposition 9.14]{BMS2} (see also \cite[Theorem 10.5]{ABGHLM}). However, we will need to study $\psi^p$ in the $p$-complete setting.

\begin{remark}[Lurie, c.f.\ the paragraph preceding {\cite[\sec IV.3]{NikolausScholze}}]
  The composite $RN$ is a map of $\T$-equivariant $\E_\infty$-rings, so it factors through the $C_p$ orbits of $\THH$ in the category of $\E_\infty$-rings, which we denote by $\THH_{\tnsr hC_p}$. The induced map $\THH_{\tnsr hC_p} \to \THH^{\Phi C_p}$ is an equivalence. If we identify the domain with $\THH$ and further map the codomain to $\THH^{tC_p}$, we obtain the Nikolaus-Scholze Frobenius.
  \[\xymatrix{
      \THH \ar[r]^-N \ar[d]_-{\psi^p} & \TR^1 \ar[r]^-R & \THH^{\Phi C_p} \ar[d]^-{\varphi_{\mathrm{NS}}}\\
      \THH_{\tnsr hC_p} \ar@{..>}[urr]_-\simeq \ar[rr]_-{\varphi_{\mathrm{NS}}} && \THH^{tC_p}
  }\]
  It follows that $RN\colon \THH \to \THH^{\Phi C_p} \simeq \THH_{\tnsr hC_p}$ is given by the $p$th Adams operation $\psi^p$.
\end{remark}

%We study Adams operations in Appendix \ref{sec:adams}, and show that $\psi^p$ acts as multiplication by $p^i$ in weight $i$. Now we will study Norm maps in gradually increasing generality, explaining the various difficulties one-by-one.

\begin{lemma}
  \label{lem:adams-perfd}
  Let $R$ be a perfectoid ring. Then $\psi^p$ acts as multiplication by $p^i$ on $\pi_{2i}\THH(R;\Z_p)$.
\end{lemma}
\begin{proof}
  The proof of \cite[Proposition 6.4.12]{RaksitDDR} goes through to show that $\psi^p$ acts as $p^i$ on $\gr^i_\HKR \HH_{\fil}(B/A)$ for any map $A\to B$ of derived commutative rings (the $\HC^-$ and $\HP$ parts of his proof do not go through). Applying this to $\HH(R;\Z_p)$, we learn that $\psi^p$ acts as $p^i$ on $\HH_{2i}(R;\Z_p)$. But the natural map $\THH_2(R;\Z_p) \to \HH_2(R;\Z_p)$ is an isomorphism by \cite[Proposition IV.4.2]{NikolausScholze}, so $\psi^p$ acts as $p^i$ on $\THH_{2i}(R;\Z_p) = \Sym^i \pi_2 \THH(R;\Z_p) = R\tnsr \xi^i A$.
\end{proof}

Let $S$ be $p$-completely smooth over any perfectoid ring $R$. In this case, recall from \cite[Corollary 6.9]{BMS2} (but essentially due to Hesselholt \cite{LarsPTyp}) that
\[ \pi_* \THH(S;\Z_p) = \Omega^*_{S/R} \tnsr_R \pi_*\THH(R;\Z_p). \]

\iffalse
\begin{lemma}
  \label{lem:adams-qrsp}
  Let $R$ be a quasiregular semiperfectoid ring. Then $\psi^p$ acts as multiplication by $p^i$ on $\pi_{2i}\THH(S;\Z_p)$.
\end{lemma}
\begin{proof}
  Choose a perfectoid ring $R$ mapping to $S$. As in the proofs of \cite[Proposition 7.10]{Prismatic} and \cite[Lemma 13.2]{Prismatic}, we can reduce to the case that $S=R\<X^{1/p^\infty}\>/(X)$. By \cite[Proposition 2.4.9]{APC}, the map of prisms $\prism_S \to \prism_S \amalg (\Z_p[q^{1/p^\infty}]^\land_{(p,q-1)},(p)_q)$ is faithfully flat, so we can assume that $R=\Z_p[\zeta_{p^\infty}]^\land_p$. In this case, $\prism_S$ is explicitly identified as a $q$-divided power envelope in \cite[\sec12.2]{Prismatic}. In particular, its Frobenius is injective, so it injects into its direct-limit perfection \cite[Lemma 3.9]{Prismatic}. Now the claim follows from Lemma \ref{lem:adams-perfd}.
\end{proof}

\todo[inline]{now extend to the proper Tate case}
\fi

\begin{lemma}
  \label{lem:adams-smooth}
  Let $S$ be $p$-completely smooth over any perfectoid ring $R$. Then $\psi^p$ acts as multiplication by $p^{i+j}$ on $\Omega^j_{S/R} \tnsr \xi^i A \subset \THH_{2i+j}(S;\Z_p)$.
\end{lemma}
\begin{proof}
  Given Lemma \ref{lem:adams-perfd}, we just need to show that $\psi^p$ acts as $p^j$ on $\Omega^j_{S/R}$, which follows by applying \cite[Proposition 6.4.12]{RaksitDDR} to $\HH(S/R;\Z_p)$.
\end{proof}

\iffalse
\subsection{Preparations}
\label{sub:prep}

\subsubsection{Induction}
Let us begin by recording how induction works.

\begin{proposition}
  \label{prop:ind-int}
  Induction of representations is given by
  \[ \ind^{n+k}_n(\lambda^j) = \sum_{r=0}^{p^k-1} \lambda^{rp^n+j}. \]
\end{proposition}
\begin{proof}
  This follows from $\res^G_H \ind^G_H V = |G/H|V$ and $(\ind^G_H V)^G = V^H$.
\end{proof}

\begin{corollary}
  \label{prop:ind-dim}
  In $p$-typical dimensional notation, induction of $C_{p^n}$-representations to $C_{p^{n+k}}$ is given by
  \[ \ind^{n+k}_n(d_0,\dotsc,d_{n-1};d_\infty) = (p^kd_0,\dotsc,p^kd_{n-1},p^kd_\infty, p^{k-1}d_\infty,\dotsc,pd_\infty;d_\infty). \]
\end{corollary}

\begin{corollary}
  \label{prop:ind-dim-poly}
  In $p$-typical dimensional notation, induction of $\Z/p^n$-polyrepresentations to $\Z/p^{n+k}$ is
  \[ \ind^{n+k}_{n}(\m d_0, \dotsc, \m d_n) = (\m d_0, p\m d_0, \dotsc, p^k d_0, p^k\m d_1, \dotsc, p^k\m d_n) \]
\end{corollary}

\subsubsection{Prismatic factorials}

\subsubsection{Refraction}

Our first goal is to find an analogue of refraction (the phenomenon that $p\in I + \Nyg^{\ge p}$ as opposed to just $I + \Nyg^{\ge 1}$). This was discovered by Mao \cite[Lemma 2.1]{MaoLogarithm}, but we will need a richer version. Adapting his proof, we get:

\begin{lemma}
  Let $R$ be a perfectoid ring, and let $\xi$ be a generator of $\Nyg^{\ge 1}\Ainf(R)$. Then for any $0\le i\le p$, we have
  \[ \phi^2(\xi) = \xi^{pi} \phi(\xi)^{p-i} + p\dfrac{\delta(\xi^i \phi(\xi))}{\phi(\xi)^i}. \]
  Moreover, $\dfrac{\delta(\xi^i \phi(\xi))}{\phi(\xi)^i}$ is a unit. In fact, we have
  \[
    \dfrac{\delta(\xi^i \phi(\xi))}{\phi(\xi)^i}
    \equiv
    \delta(\xi)^\phi
    \bmod \phi(\xi)^{p-i}
  \]
\end{lemma}
\begin{proof}
  The equation follows by writing $\phi^2(\xi) = \frac{\phi(\xi^i \phi(\xi))}{\phi(\xi)^i}$. On the other hand, we have
  \begin{align*}
    \phi^2(\xi)
      &= (\xi^p + p\delta(\xi))^i \phi(\xi)^{p-i} + p\delta(\xi)^\phi\\
      &= \xi^{pi} \phi(\xi)^{p-i} + p\left[\delta(\xi)^\phi + \phi(\xi)^{p-i}\fade{\sum_{j=0}^{i-1} \binom ij \xi^{pj} p^{(i-1)-j} \delta(\xi)^{i-j}}\right]
  \end{align*}
  \todo[inline]{justify dividing by $p$}
\end{proof}
\fi

\subsection{de Rham-Witt norms}
\label{sub:dRW-norm}

In this subsection, we describe the Norm maps that the de Rham-Witt complex inherits from $\TR$.

First, we record how induction works.

\begin{proposition}
  In $p$-typical notation, induction of polyrepresentations is given by
  \[ \ind_n^{n+k}(\m d_0,\dotsc, \m d_n) = (\m d_0, p\m d_0, \dotsc, p^k \m d_0, p^k \m d_1, \dotsc, p^k \m d_n). \]
  Integrally, we have
  \[ \ind^{\Z/mn}_{\Z/n}(\lambda^j) = \sum_{i=0}^{m-1} \lambda^{in+j} \]
\end{proposition}
\begin{proof}
  These follow from $\res^G_H \ind^G_H \alpha = |G/H|\alpha$ and $(\ind^G_H\alpha)^G = \alpha^H$.
\end{proof}

It follows that there are Norm maps
\[ \dRW[n]^j_{S/R} \mono \TR^n_j(S;\Z_p) \morph^N \TR^{n+1}_{\ind j}(S;\Z_p) \hookleftarrow \dRW[n+1]^j_{S/R} \fade{{}\tnsr({}\cdots)} + \Fil^1\dRW[n+1]^{pj}_{S/R} \fade{{}\tnsr({}\cdots)}\]
where $({}\cdots)$ is some complicated invertible module which we ignore for the moment. There are other terms in $\TR^{n+1}_{\ind j}$, but only these two are relevant to the identities $FN\omega = \omega^p$ and $N \equiv \psi^p \bmod \Fil^1$.

As explained in the Introduction, the operation
\begin{align*}
  N(\omega)
    &\defeq (p-V(1))\frac{\psi^p(\omega)}p + V\frac{\omega^p}p\\
    &= \psi^p(\omega) - V\underbrace{\frac{F\psi^p\omega - \omega^p}p}_{{}\defeq \delta(\omega)}
\end{align*}
satisfies these identities. Beware that $V$ here is the $\RO(G)$-graded Verschiebung
\[ V\colon \TR^n_{pj}(S;\Z_p) \to \TR^{n+1}_{\ind(j)}(S;\Z_p). \]
Before identifying the twists $({}\cdots)$, let us explain why these identities suffice to pin down the Norm.

\begin{lemma}
  \label{lem:norm-ghost}
  Suppose $N\colon \TR^n_j(S;\Z_p) \to \TR^{n+1}_{\ind(j)}(S;\Z_p)$ is a natural transformation of functors
  \[ \{\text{\normalfont rings $p$-completely smooth over some perfectoid ring}\} \to \{\text{\normalfont sets}\}.\]
  satisfying $FN\omega = \omega^p$ and $N \equiv \psi^p \bmod \Fil^1$. Then $N$ coincides with the Hill-Hopkins-Ravenel norm. 
\end{lemma}
\begin{proof}
  Let $S$ be $p$-completely smooth over the perfectoid ring $R$. When $R$ is $p$-torsionfree, so is $S$, and thus the maps $F$ and $R^{n+1}$ are jointly injective, so the identities uniquely determine $N$.

  For a general perfectoid ring $R$, there exists a surjection $R^+ \epi R$ from a $p$-torsionfree perfectoid ring $R^+$ by \cite[Proposition 2.12]{CSPurity}. The $p$-completed cotangent complex $\L_{\Spf R/\Spf R^+} = (\L_{R/R^+})^\land_p$ vanishes by \cite[Lemma 3.14]{BMS1}, so by deformation theory $S$ descends to a $p$-completely smooth $R^+$-algebra $S^+$. By \cite[Proposition 10.14]{BMS1}, the natural map
  \[ \dRW[n]^*_{S^+/R^+} \mathbin{\hat\tnsr}_{\rmW_n(R^+)} \rmW_n(R) \to \dRW[n]^*_{S/R} \]
  is an isomorphism. Since $\rmW_n(-)$ preserves surjections, $\rmW_n(R^+) \to \rmW_n(R)$ is surjective, hence so is $\dRW[n]^*_{S^+/R^+} \to \dRW[n]^*_{S/R}$. This shows that the Norm for $S^+$ uniquely determines the Norm for $S$.
\end{proof}

In the following theorem, it is not obvious that the various forms of the codomain are equivalent; this will be explained in the proof. Recall the prismatic factorial notation from Definition \ref{dfn:prismatic-factorial}, and also that we write $n? \defeq (n-1)!$ for the Gamma function.

\begin{theorem}[Norms for the de Rham-Witt complex]
  \label{thm:norms-dRW}
  Let $R$ be any perfectoid ring, and $S$ a $p$-completely smooth $R$-algebra. When $p>2$ and $j\in\N$, there are Norm maps
  \begin{align*}
    N^k = N^{n+k}_n\colon
    \dRW[n]^{j}_{S/R}
    &\morph
    \sum_{r=0}^k
    \Fil^r \dRW[n+k]^{jp^r}_{S/R} \tnsr \left(\frac{\phi^{-k}\left(\xi_{n+1}^{h_1} (h_1)_\prism!\right)}{\xi_{n+k+1}^{(p^r-1)/2}}A\right)^{\tnsr j}\\
    &\hspace{.6em}=\fade{\sum_{r=0}^k
    \Fil^r \dRW[n+k]^{jp^r}_{S/R} \tnsr{}} \frac{\phi^{-k}\left(\xi_{n+1}^{h_j} (h_j)_\prism!\right)}{\xi_{n+k+1}^{(p^rj-1)/2}}A \qquad\text{\normalfont for $j$ odd}
    \intertext{where $h_m\defeq\frac{p^k m-1}2$, and $\sum$ is $\bigoplus$ for $j>0$. For even degrees, this can equivalently be written}
    \fade{N^k = N^{n+k}_n\colon}
    \dRW[n]^{2j}_{S/R}
    &\morph
    \sum_{r=0}^k
    \Fil^r \dRW[n+k]^{2jp^r}_{S/R} \tnsr \left(\frac{\phi^{-k}\left(\xi_{n+1}^{p^k-1} (p^k)_\prism?\right)}{\xi_{n+k+1}^{p^r-1}}A\right)^{\tnsr j}\\
    &\hspace{.6em}=\fade{
    \sum_{r=0}^k
  \Fil^r \dRW[n+k]^{2jp^r}_{S/R} \tnsr{}} \frac{\phi^{-k}\left(\xi_{n+1}^{p^kj-1} (p^k j)_\prism?\right)}{\xi_{n+k+1}^{p^rj-1}}A\\
  \intertext{and this form also works for $p=2$.}
  \intertext{When $R=k$ is a perfect $\F_p$-algebra and $j>0$, these reduce to an operation}
    % general form
    N^k \colon
    \dRW[n]^{j}_S
    &\morph
    \Fil^k \dRW[n+k]^{jp^k}_S \tnsr \left(\frac{h_1!}{p^{kh_1}}W\right)^{\tnsr j}\\
    % odd form
    &\hspace{.6em}=\fade{
    \Fil^k \dRW[n+k]^{jp^k}_S \tnsr{}} \frac{h_j!}{p^{kh_j}}W \qquad \text{\normalfont for $j$ odd}\\
    % even form 1
    \fade{N^k \colon}
    \dRW[n]^{2j}_S
    &\morph
    \Fil^k \dRW[n+k]^{2jp^k}_S \tnsr \left(\frac{p^k?}{p^{k(p^k-1)}}W\right)^{\tnsr j}\\
    % even form 2
    &\hspace{.6em}=\fade{
    \Fil^k \dRW[n+k]^{2jp^k}_S \tnsr{}} \frac{(p^kj)?}{p^{k(p^kj-1)}}W\\
    \intertext{which also exists for $n=\infty$.}
    \intertext{When $R=\O_C$ is spherically complete, they assemble to an infinite Norm map}
    % general form
    N^k \colon
    \dRW^{j}_{S/\O_C}
    &\morph
    \sum_{r=0}^k
    \Fil^r \dRW^{jp^r}_{S/\O_C} \tnsr \left(\frac{\phi^{-k}\left(\mu^{h_1} (h_1)_q!\right)}{\mu^{(p^r-1)/2}}\Ainf(\O_C)\right)^{\tnsr j}\\
    % odd form
    &\hspace{.6em}=\fade{\sum_{r=0}^k
    \Fil^r \dRW^{jp^r}_{S/\O_C} \tnsr{}} \frac{\phi^{-k}\left(\mu^{h_j} (h_j)_q!\right)}{\mu^{(p^rj-1)/2}}\Ainf(\O_C) \qquad\text{\normalfont for $j$ odd}\\
    % even form 1
    \fade{N^k \colon}
    \dRW^{2j}_{S/\O_C}
    &\morph
    \sum_{r=0}^k
    \Fil^r \dRW^{2jp^r}_{S/\O_C} \tnsr \left(\frac{\phi^{-k}\left(\mu^{p^k-1} (p^k)_q?\right)}{\mu^{p^r-1}}\Ainf(\O_C)\right)^{\tnsr j}\\
    % even form 2
    &\hspace{.6em}=\fade{
    \sum_{r=0}^k
  \Fil^r \dRW^{2jp^r}_{S/\O_C} \tnsr{}} \frac{\phi^{-k}\left(\mu^{p^kj-1} (p^k j)_q?\right)}{\mu^{p^rj-1}}\Ainf(\O_C)
  \end{align*}

  These maps satisfy the following properties:
  \begin{itemize}
    \item when $j=0$, $N(f)=f - V\delta(f)$ is the Norm map of Witt vectors \cite{AngeltveitNorm};

    \item when $j>0$, $N$ is given by 
      \begin{align*}
        N(\omega)
          &= (p-V(1))\frac{\psi^p(\omega)}p + V\frac{\omega^p}p
          % &= (p-V(1))p^{2j-1}\omega + V\frac{\omega^p}p
      \end{align*}
      or more generally
      \begin{align*}
        N^k(\omega)
          &= \sum_{r=0}^k V^r(p^{k-r} - V{\lf p^{k-r-1}\rf})\frac{\psi^{p^{k-r}}(\omega)^{p^{r}}}{p^k}
          % &= \sum_{r=0}^k V^r(p^{k-r} - V{\lf p^{k-r-1}\rf}) p^{2(k-r)jp^r - k}\omega^{p^r}
      \end{align*}

    \item when $R=k$ is a perfect $\F_p$-algebra, the preceding formula simplifies to
      \[ N^k(\omega) = V^k \frac{\omega^{p^k}}{p^k} \]
      for $|\omega|>0$, and in particular $N(\omega)=0$ whenever $\omega$ is of odd degree;

    \item $N$ satisfies the following identities for all $\omega,\,\tau$:
      \begin{align*}
        FN\omega &= \omega^p\\
        NV\omega &= p^{p-2}V^2\omega^p\\
        NR\omega &= RN\omega\\
        N(\omega\tau) &= N(\omega)N(\tau)\\
        N(\omega+\tau) &= N(\omega)+N(\tau)+Vs_1(\omega, \tau)\\
                       &\phantom{={}}\text{ where } (a+b)^p = a^p + b^p + ps_1(a, b)
      \end{align*}
  \end{itemize}
\end{theorem}

\begin{proof}
  We have already explained why the formula for $N$ must hold, and the final list of identities is true for any Tambara functor. It remains to identify the twists in the codomain of $N$. First suppose that $p$ is odd and $j=1$. We have
  \[ \ind^{n+k}_n(1_\R) = 1_\R + \lambda^{p^n} + \dotsb + \lambda^{h_1p^n} \]
  so that $\ind^{n+k}_n(1) - 1$ corresponds in prismatic notation to the twist by
  \begin{align*}
    (p^{n+1})_{\prism^{1/p}} \tnsr \dotsb \tnsr \left(h_1 p^{n+1}\right)_{\prism^{1/p}}
    &=
    \xi^{h_1} I_n^{h_1} \tnsr (h_1)_{\prism^{p^n}}!
  \end{align*}
  To switch to polygonic notation, we apply $\phi^{-(n+k)}$, turning this into
  \[ \phi^{-k}\left(\xi_{n+1}^{h_1} \tnsr (h_1)_\prism!\right) \]
  This is the twist for the $\dRW[n+k]^1_{S/R}$ term in the codomain of $N^k\colon \dRW[n]^1_{S/R} \to ({}\cdots)$. Each additional $\Fil^r \dRW[n+k]^{p^r}_{S/R}$ term moves $\frac{p^r-1}2\lambda^0$ from the topological part to the de Rham-Witt part, which has the effect of dividing by $\xi_{n+k+1}^{(p^r-1)/2}$.

 That is the analysis for $j=1$; for $j>1$, we can just raise everything to the $j^\th$ power, which yields the first form. But for $j$ odd, we could also write
 \[ \ind^{n+k}_n(j_\R) = 1_\R + \lambda^{p^n} + \dotsb + \lambda^{h_j p^n} \]
 which yields a twist of
 \[\xymatrix@C=4em{
     \xi^{h_j} I_n^{h_j} \tnsr (h_j)_{\prism^{p^n}}!
    \ar@{~>}[r]^-{\phi^{-(n+k)}} &
    \phi^{-k}\left(\xi_{n+1}^{h_j} \tnsr (h_j)_{\prism}!\right)
 }\]
 Then we need to subtract $\frac{p^rj-1}2 \lambda^0$ to land in $\Fil^r\dRW[n+k]^{jp^r}_{S/R}$, which yields the second form.

 Similarly, for even degrees $2j$ we can write
 \begin{align*}
   \ind^{n+k}_n(j\lambda^0)
    &= j\ind^{n+k}_n(\lambda^0)\\
    &= j\left(\lambda^0 + \lambda^{p^n} + \dotsb + \lambda^{(p^k-1)p^n}\right)
    \intertext{or}
    \ind^{n+k}_n(j\lambda^0)
    &= \lambda^0 + \lambda^{p^n} + \dotsb + \lambda^{(p^kj-1)p^n}
\end{align*}
which give the third and fourth forms respectively.

It is not obvious that the various forms of the twist are equivalent. In fact, they are not if the symbols are interpreted as $C_{p^\infty}$-Mackey functors. But they are equivalent when restricted to $C_{p^{n+k}}$, which is all that matters. For a simple example, $((\frac{5-1}2)!)^2=4$ is equal to $(5-1)!=24$ in $\Z/5$ but not in $\Z/25$; beware that $(2)_q$ is not equal to $(3)_q (4)_q$ in $\Z[q]/(5)_q$. Evidently many other formulations of the twist are possible.

  The claims about the Norm simplifying in characteristic $p$ are nontrivial, and proved in Corollary \ref{cor:norms-Fp} below. But we can identify the codomain of the $n=\infty$ Norms over $k$ and over $\O_C$. Over $\O_C$, we just replace every instance of $\xi_\bullet$ by $\mu$. In characteristic $p$, we have
  \[
    \frac{\phi^{-(k+1)}\left(\xi_n^{h_1} (h_1)_\prism!\right)}{\xi_{n+k+1}^{(p^r-1)/2}}
    =
    \frac{p^{nh_1} h_1!}{p^{(n+k+1)(p^r-1)/2}}
  \]
  and
  \[ n\frac{p^k-1}2 - (n+k+1)\frac{p^r-1}2 = n\frac{p^k - p^r}2 - (k+1)\frac{p^r-1}2 \]
  which goes to $\infty$ for $r<k$. Thus only the $r=k$ summand survives the limit in characteristic $p$.
\end{proof}

Next we want to justify the claim that $N(\omega)=V\frac{\omega^p}p$ in characteristic $p$. Intuitively this is true because $p-V(1)=0$ in $\rmW(k)$; however, we must be careful since some of this factor gets absorbed into the twist. Instead, as in Lemma \ref{lem:norm-ghost}, we must use $p$-torsionfree perfectoid rings to control the Norm in characteristic $p$.

For that matter, even in mixed characteristic, there is an implicit claim that $p-V(1)$ is divisible by the ideal twist. The resolution is a subtle property that has come up repeatedly when considering the Norm.

\begin{remark}[Refraction]
  \label{rmk:snell}
  There are two familiar instances of the optical phenomenon of \emph{refraction}\footnote{Previously the author referred to the first as diffraction and the second as refraction, but was informed that this is incorrect.}, with different corresponding mathematical incarnations. First, a prism breaks up white light into constituent wavelengths; according to Bhatt, the famous depiction of this on the cover of \cite{DarkSideMoon} was the inspiration for the name \emph{prismatic} cohomology, specifically due to the axiom $p \in I + \phi(I)$. In terms of Mackey data, the Tambara functor $\m{(A, I)}$ constructed in \cite{PrismsTambara1} has the feature that $\ker F^{n+k}_n$ and $\im V^{n+k}_n$ are \emph{annihilators} of each other; for a generic Green functor, all one can say is that $(\ker F^{n+k}_n)(\im V^{n+k}_n) = 0$ (by Frobenius reciprocity). We refer to this incarnation of refraction as \emph{Floyd refraction}, and say that the Green functor $\m{(A, I)}$ is \emph{Floyd refractive}.

  The second instance is that objects placed in water seem larger than they actually are; this distortion is described by Snell's law. Mathematically, this corresponds to being able to strengthen the prism condition to $p \in I^p + \phi(I)$; a generic Green functor always has $p - V(1) \in \ker F$, but $\res^{C_{p^\infty}}_{C_p} \m{(A, I)}$ obeys the scaled condition $p-V(1) \in (\ker F)^p$. We refer to this phenomenon as \emph{Snell refraction}.

  We highlighted this phenomenon in \cite{PrismsTambara1}, but there we did not know how to generalize Snell refraction beyond $C_p$. Mao discovered the correct generalization when using the norm to construct the (refined) prismatic logarithm \cite[Lemma 2.1]{MaoLogarithm}. His formulation is that for a transversal prism $(A, I)$,
  \[ \frac{\phi^n(I)}{I \cdot I^{p-1}_n} \to \frac{A}{I \cdot I_r^{p-1} } \]
  is a monomorphism with image $(p)$. Examining the proof, we find that in fact
  \[ p = \phi^n(\~\xi)\fade u + \~\xi \~\xi_n^{p-1}\fade v \]
  for units $u$ and $v$. Applying $\phi^{-(n+1)}$, we get
  \begin{equation}
    \label{eq:snell-polygonic}
    p = \xi \fade u + \phi^{-1}(\xi_n)^{p-1} \phi^{-n}(\xi) \fade v.
  \end{equation}
  This is the form we will use.
\end{remark}

\begin{corollary}
  \label{cor:norms-Fp}
  In characteristic $p$, for $|\omega|>0$ the Norm is given simply by $N^k(\omega)=V^k\frac{\omega^{p^k}}{p^k}$. Furthermore, $N(\omega)=0$ for $|\omega|$ odd.
\end{corollary}
\begin{proof}
  Let $R$ be a $p$-torsionfree perfectoid ring, and let $j>0$. First suppose that $p$ is odd, and let $h=\frac{p-1}2$ as above. Then the Norm goes
  \[
    N\colon \dRW[n]^j_{S/R} \to
    \dRW[n+1]^j_{S/R} \tnsr \phi^{-1}\left(\xi_{n+1}^h (h)_\prism!\right)^{\tnsr j}
    \fade{{}\oplus
    \Fil^1\dRW[n+1]^{pj}_{S/R} \tnsr ({}\cdots)}
  \]
  We can ignore the factor $\phi^{-1}((h)_\prism!)^j$ which is a unit. Now the expression ``$(p-V(1))\frac{\psi^p(\omega)}p$'' means an expression that goes to $0$ under $F$ and to $\psi^p(\omega)$ under $R^{n+1}$. Using Snell refraction \eqref{eq:snell-polygonic},
  \[
    \phi^{-1}(\xi_{n+1})^{hj} \phi^{-(n+1)}(\xi)^j \omega \tnsr \phi^{-1}(\xi_{n+1})^{hj} \approx p^j \omega \bmod \xi
  \]
  is (up to units) such an expression, since the $\phi^{-1}(\xi_{n+1})$ term is killed by $F$. Specializing to characteristic $p$, the coefficient of $\omega$ in the first tensor factor becomes
  \[ p^{(n+1)hj+j} \ge p^{(n+1)h+1} \ge p^{n+2} = 0 \in \rmW_{n+1}(k). \]

  When $p=2$ and $j$ is even, $hj$ still makes sense, and the above argument still goes through, except the last line becomes $p^{(n+1)hj+j} \ge p^{(n+1)+2} = 0 \in \rmW_{n+1}(k)$. We do not make any claims about odd-degree Norms when $p=2$. 

  To handle iterates of the Norm, one can either generalize the above argument, or just induct using the identity $NV\omega = p^{p-2} V^2 \omega^p$. To see that $N(\omega)=0$ for odd $|\omega|$, it suffices to show that $\omega \mapsto \frac{\omega^{p^k}}{p^k}$ is zero for odd $|\omega|$. This is easy to see by reduction to the $p$-torsionfree case.
\end{proof}

